\section{Conditioning the finite realization}
\label{sec:conditioning}

The approximation theorem separates existence of convergence from the
size of a useful error budget. Its constants are explicit, but an
absolute product envelope centered at the origin can be expensive when
the observed pair has large real part. We first identify that loss,
then derive a centered estimate for the \emph{same} initialized
polynomial. The complete Gaussian path, the calibration and the
eventual transport consumer are all retained
\cite{BakerCM50,BakerCM51}. The fixed numerical instance is deferred to
Appendix~\ref{supp:conditioning}; no numerical constant below is needed
for the approximation theorem.

\subsection{A limitation of the global absolute allowance}

Write $h_0=H(t_0,0)>0$, and let $h_H>0$ be a common floor on the
counted contours. Each circle meets the real axis, where positivity of
the source gives $|H(t_0,x)|\le h_0$. Thus $h_H\le h_0$.
The Gaussian constant of Section~\ref{sec:common-time}, with
$\lambda=[8(1+\Delta)]^{-1}$, contains the positive contribution
\[
 2M_\lambda R^2\sqrt{1+\Delta}
       \exp\!\left(\frac{R^2}{4(1+\Delta)}\right),
 \qquad M_\lambda\ge h_0.
\]
Here the origin-centered disk of radius $R$ contains the contours.
When the omitted tail is replaced by the count allowance
\[
 \overline E_N=\frac{2K\{\log(2+N)+1\}}{N},
 \qquad K=5C_{\mathrm u}\ge5,
\]
the resulting global allowance $\varepsilon_{\rm glob}$ satisfies
\begin{equation}
 \frac{\varepsilon_{\rm glob}}{h_H}
 \ge \frac{20R^2\sqrt{1+\Delta}}{N}
          \exp\!\left(\frac{R^2}{4(1+\Delta)}\right).
 \label{cond:global-screen}
\end{equation}
Consequently this particular allowance can meet
$\varepsilon_{\rm glob}\le h_H/2$ only if
\[
 N\ge40R^2\sqrt{1+\Delta}
          \exp\!\left(\frac{R^2}{4(1+\Delta)}\right).
\]
This lower-bounds a cutoff required by the \emph{displayed sufficient
allowance}; it does not lower-bound the actual approximation error or
the necessary complexity of a finite realization. Multiplying both
fields by one fixed positive scalar multiplies their error and floor
equally and does not change this ratio.

\subsection{A fixed scalar calibration}

We now specialize the initialization to $t_0=0$ and put
$\widetilde H=gH$ for a fixed $g>0$. Let the observation disk be
$|z-x_*|\le d_0$, with $x_*>0$, and choose
$R\ge\max(1,x_*+d_0)$. Assume $0<\Delta\le1/2$ and put
$v=2\sqrt{\Delta}$. For $N>2R+1$, define
\begin{equation}
 F_N(w)=\prod_{|\rho_\nu|>N}(1-w^2/\rho_\nu^2),\quad
 A_N=F_N(x_*),\quad
 L_N=\frac{N/2-R}{v},\quad E=E_N.
 \label{cond:calibration}
\end{equation}
The product and tail use one representative per opposite occurrence,
as in Section~\ref{sec:common-time}; all ties and conjugates remain.

\begin{lemma}[Calibrated omitted-product estimate]\label{cond:ratio}
The scalar $A_N$ is real and satisfies
\begin{equation}
 0<A_N\le1,\qquad
 |1-A_N|\le-\log A_N\le\frac43x_*^2E.
 \label{cond:scalar}
\end{equation}
For $0\le t\le\Delta$, $|z-x_*|\le d_0$ and $|r|\le L_N$,
put $w=z+2\ii\sqrt t\,r$. Then
\begin{align}
 |\log F_N(w)-\log F_N(x_*)|&\le E p(|r|),\notag\\
 p(u)&=c_0+c_1u+c_2u^2,\notag\\
 c_0&=\tfrac43d_0(2x_*+d_0),\quad
 c_1=\tfrac83(x_*+d_0)v,\quad c_2=\tfrac43v^2.
 \label{cond:local-ratio}
\end{align}
\end{lemma}

\begin{proof}
Both $w$ and $x_*$ lie in $|w|\le N/2$. Along the straight segment
between their squares every omitted factor has denominator at least
$3/4$ in modulus. The convergent logarithmic series therefore gives
\[
 |\log F_N(w)-\log F_N(x_*)|
 \le \tfrac43E|w^2-x_*^2|
 \le\tfrac43E(d_0+v|r|)(2x_*+d_0+v|r|).
\]
For an omitted zero the strip implies
$\Re(\rho_\nu^2)=|\rho_\nu|^2-2(\Im\rho_\nu)^2>N^2-2$.
In particular $2\Re(\rho_\nu^2)>x_*^2$, so
$|1-x_*^2/\rho_\nu^2|<1$. Conjugate factors pair positively, and
the real factors are positive because $|\rho_\nu|>N>x_*$.
The absolutely convergent real logarithm proves $A_N>0$, while
termwise moduli give $A_N\le1$. Finally,
\[
 -\log A_N\le\sum_\nu
       \frac{x_*^2/|\rho_\nu|^2}{1-x_*^2/N^2}
       \le\tfrac43x_*^2E .
\]
The inequality $1-A_N\le-\log A_N$ completes the proof.
\end{proof}

Calibration changes no zero, rational moment or heat evolution:
the comparison polynomial is $\widetilde P_N=gA_NP_N$, with one
constant multiplier for all times and all spatial derivatives. At the
initial time the exact identity
\begin{equation}
 \widetilde H(0,w)-gA_NP_N(0,w)
 =\widetilde H(0,w)
    \{1-\exp(\log F_N(x_*)-\log F_N(w))\}
 \label{cond:exact-ratio}
\end{equation}
does not divide by a retained factor. It remains valid at retained
zeros. In particular separate calibrations for the field and derivative
would not be legitimate inputs to a logarithmic-derivative quotient.

\subsection{Evaluating the complete Gaussian path}

Assume the following independently supplied source bounds, for some
$u_0>0$ and positive $B_c,C,\alpha$:
\begin{align}
 |\widetilde H(0,w)|&\le B_c
       &&(|w-x_*|\le d_0+vu_0),\notag\\
 |\e^{-r^2}\widetilde H(0,z+2\ii\sqrt t\,r)|
     &\le C\e^{-\alpha r^2+\beta|r|}
       &&(|z-x_*|\le d_0,\ 0\le t\le\Delta,\ r\in\R).
 \label{cond:source-majorants}
\end{align}
Here $\beta\ge0$, and we require $L_N>u_0$.
Choose a fixed $E_*>0$ such that
$1-E_*c_2>0$ and $\alpha-E_*c_2>0$.
Define the explicit Gaussian integrals
\begin{align}
 J(A,B;l)
 &=\frac{\sqrt\pi}{2\sqrt A}
       \e^{B^2/(4A)}
       \operatorname{erfc}\!\left(\sqrt A\,l-\frac B{2\sqrt A}\right),
       \qquad A>0,\notag\\
 I_0&=J,\qquad
 I_1(A,B;l)=\frac{\e^{-Al^2+Bl}+BJ(A,B;l)}{2A},\notag\\
 I_2(A,B;l)&=\frac{l\e^{-Al^2+Bl}+J(A,B;l)+BI_1(A,B;l)}{2A}.
 \label{cond:gaussian-integrals}
\end{align}
Completing the square gives $J=\int_l^\infty\e^{-Au^2+Bu}\dd u$.
Integration by parts then identifies $I_j$ with the same integral
weighted by $u^j$, for $j=1,2$.

\begin{proposition}[Centered complete-path bound]\label{cond:path}
Let $\mathcal G(s)$ be the absolute-product envelope from
Section~\ref{sec:common-time}, and choose $M_\lambda$ there so that
$\mathcal G(s)\le M_\lambda\e^{\lambda s^2}$, where
$\lambda=[8(1+\Delta)]^{-1}$ and
$b=1-8\lambda\Delta=(1+\Delta)^{-1}$.
For $0\le E\le E_*$,
\begin{equation}
 \sup_{\substack{0\le t\le\Delta\\|z-x_*|\le d_0}}
 |\widetilde H(t,z)-gA_NP_N(t,z)|
 \le \mathcal A_{\rm cent}(E)+E K_{\rm far}(N),
 \label{cond:complete-bound}
\end{equation}
where
\begin{align}
 \mathcal A_{\rm cent}(E)
 &=\frac{2B_c}{\sqrt\pi}
      \{\e^{Ec_0}J(1-Ec_2,Ec_1;0)-J(1,0;0)\}\notag\\
 &\quad+\frac{2C}{\sqrt\pi}
      \{\e^{Ec_0}J(\alpha-Ec_2,\beta+Ec_1;u_0)
                        -J(\alpha,\beta;u_0)\},\notag\\
 K_{\rm far}(N)
 &=gM_\lambda\e^{2\lambda R^2}
   \{(2R^2+\tfrac43x_*^2)J_0(b,L_N)+2v^2J_2(b,L_N)\},\notag\\
 J_0(b,L)&=b^{-1/2}\operatorname{erfc}(\sqrt b\,L),\notag\\
 J_2(b,L)&=\frac{L\e^{-bL^2}}{\sqrt\pi b}
            +\frac{\operatorname{erfc}(\sqrt b\,L)}{2b^{3/2}} .
 \label{cond:complete-constants}
\end{align}
The near-path term is at most $E K_{\rm cent}$, with
\begin{align}
 K_{\rm cent}
 =\frac{2\e^{E_*c_0}}{\sqrt\pi}\bigg\{&
 B_c\sum_{j=0}^2 c_j I_j(1-E_*c_2,E_*c_1;0)\notag\\
 &+C\sum_{j=0}^2 c_j
       I_j(\alpha-E_*c_2,\beta+E_*c_1;u_0)\bigg\}.
 \label{cond:linear-constant}
\end{align}
\end{proposition}

\begin{proof}
Use the exact imaginary Gaussian formula for both fields.
On $|r|\le u_0$ apply the first bound in
\eqref{cond:source-majorants}; on $u_0<|r|\le L_N$ apply its second.
Equations~\eqref{cond:local-ratio}--\eqref{cond:exact-ratio} bound the
respective integrands by
$B_c\e^{-r^2}(\e^{Ep(|r|)}-1)$ and
$C\e^{-\alpha r^2+\beta|r|}(\e^{Ep(|r|)}-1)$.
Extend these \emph{nonnegative majorant integrals} to $[0,\infty)$
and $[u_0,\infty)$, respectively. Their evaluated values are
$\mathcal A_{\rm cent}(E)$. Neither source bound nor the product-ratio
identity is thereby asserted outside its original domain.
The inequality $\e^{Ep}-1\le Ep\e^{E_*p}$ gives
\eqref{cond:linear-constant}.

On $|r|>L_N$ no logarithmic expansion is used. Telescoping the
uncalibrated product and then using \eqref{cond:scalar} gives
\[
 |\widetilde H(0,w)-gA_NP_N(0,w)|
 \le gE(|w|^2+\tfrac43x_*^2)\mathcal G(|w|).
\]
Since $|w|^2\le2R^2+2v^2r^2$, multiplication by $\e^{-r^2}$
and integration of $1$ and $r^2$ over the two Gaussian tails gives
exactly $E K_{\rm far}(N)$. This includes the calibration mismatch.
\end{proof}

The large radial exponential has not been discarded. It occurs only
in the tail beginning at $L_N\to\infty$. To quantify its suppression,
put
\[
 B=\log(gM_\lambda)+2\lambda R^2,\qquad
 D_f=\frac{2R^2+\frac43x_*^2+v^2/b+2v^2}{\sqrt\pi b}.
\]
For $L\ge1$, the elementary bound
$\operatorname{erfc}(\sqrt bL)\le
\e^{-bL^2}/(\sqrt{\pi b}L)$ implies
$K_{\rm far}(N)\le D_fL_N\e^{B-bL_N^2}$.
Because $b\ge2/3$, $\log L\le bL^2/2$ for $L\ge8$.
Thus the explicit sufficient test
\begin{equation}
 L_N\ge
 \max\!\left\{u_0,8,
       \sqrt{\frac2b\max(0,B+\log D_f)}\right\},
 \qquad L_N>u_0,
 \label{cond:far-test}
\end{equation}
ensures $K_{\rm far}(N)\le1$. Fixed supplied inputs therefore give
a complete $(K_{\rm cent}+1)E_N$ allowance eventually in $N$.
Its main coefficient no longer contains the origin-height quadratic
exponential, but its absolute amplitude cost may remain unfavorable.

On a padded observation disk, analytic finite-data error is added to
\eqref{cond:complete-bound}. Cauchy's formula supplies the derivative
payment $j!\eta^{-j}\varepsilon_{\rm pad}$ at padding $\eta$.
Arbitrary pointwise evaluation roundoff is paid separately; it is not
an analytic function to which Cauchy's estimate can be applied.
If a represented scalar $\widehat A_N$ replaces $A_N$, add
$|\widehat A_N-A_N|$ times a bound for the entire evolved polynomial
and each required derivative. Exact latent calibration changes no
divisor, but uncertainty in acquired coefficients still carries the
weights of Section~\ref{sec:counted-transfer}.

\subsection{The consumer can fail after the tail is repaired}

The positive central allowance itself has the lower estimate
\begin{equation}
 \frac{B_c}{\sqrt\pi}\int_{-u_0}^{u_0}
       \e^{-r^2}(\e^{E p(|r|)}-1)\dd r
 \ge B_cc_0E\,\operatorname{erf}(u_0).
 \label{cond:core-floor}
\end{equation}
Indeed $\e^x-1\ge x$ for $x\ge0$, and $p\ge c_0$.
For $u_0\ge1$ this is greater than $(2/3)B_cc_0E$:
integrating $\e^{-r^2}\ge1-r^2$ on $[-1,1]$ and using
$\sqrt\pi<2$ gives the stated Gaussian mass.
If one inserts the \emph{upper allowance} $\overline E_N$ in this
formula and restricts $N\le N_{\max}$, the chosen central allowance
is therefore greater than
$4KB_cc_0/(3N_{\max})$. This is not a lower bound on $E_N$ or on
the signed Gaussian error.

To see how this matters to an actual consumer, let $h$ be a contour
field floor, $K_1$ an upper derivative bound, $r$ its radius and
$q_\Gamma$ a positive quadratic floor. The transfer of
Section~\ref{sec:counted-transfer} includes the nonnegative terms
\[
 b_1\ge\frac{K_1\varepsilon_0}{h(h-\varepsilon_0)}
       \ge\frac{K_1}{h^2}\varepsilon_0,\qquad
 e_M\ge\frac{2r}{q_\Gamma}b_1.
\]
When $\varepsilon_0\ge h/2$ the first contour test is already
unavailable. Otherwise the \emph{sufficient} every-subinterval
logarithmic transport budget is at least
\begin{equation}
 4\Delta e_M\ge
       \frac{8\Delta rK_1}{q_\Gamma h^2}\varepsilon_0 .
 \label{cond:consumer-floor}
\end{equation}
If the quadratic-denominator test fails, the transfer is unavailable
for that separate reason. If it passes, \eqref{cond:core-floor} and
\eqref{cond:consumer-floor} may already exceed a chosen tolerance,
even after granting zero derivative, data and complementary errors.
No precision increase cancels positive terms in this specified budget.
The conclusion is limited to this absolute-envelope/count-tail/quotient
allowance. Relative errors, correlated contour data, acquired tail
information and the actual approximation error remain different questions.
